% Part: first-order-logic
% Chapter: lindstrom
% Section: lindstrom-proof

\documentclass[../../../include/open-logic-section]{subfiles}

\begin{document}

\olfileid{mod}{lin}{prf}

\olsection{लिंडस्ट्रॉमचे प्रमेय}

\begin{lem}
\ollabel{lem:lindstrom}
$!E \in L(\Lang{L})$ आणि $\Lang{L}$ सांत आहे असे समजा; तसेच असा
$n \in \Nat$ आहे असे समजा की कोणत्याही दोन !!{structure}s
$\Struct{M}$ आणि~$\Struct{N}$ साठी, $\Struct{M} \equiv_n
\Struct{N}$ आणि $\Struct{M} \models_L !E$ असेल, तर $\Struct{N}
\models_L !E$ हेही खरे असते. मग $!E$ हे एका प्रथम-क्रम
!!{sentence} शी सममूल्य असते; म्हणजे असे प्रथम-क्रम $!D$ असते की
$\Mod(L){!E} = \Mod(L){!D}$.
\end{lem}

\begin{proof}
कोणत्याही दोन $n$-सममूल्य !!{structure}s $\Struct{M}$ आणि
$\Struct{N}$ या $!E$ ला नेमलेल्या सत्यतामूल्याबाबत सहमत असतील, असा $n$
घ्या. \olref[bas][pis]{prop:qr-finite} आठवा: तार्किक सममूल्यतेपर्यंत,
सांत !!{language} मध्ये संख्यापक-प्रतांक $n$ पेक्षा अधिक नसलेली प्रथम-क्रम
!!{sentence}s सांत इतकीच असतात. आता प्रत्येक स्थिर
!!{structure}~$\Struct{M}$ साठी, $\Struct{M}$ मध्ये सत्य आणि
$\QuantRank{!E} \le n$ असलेल्या सर्व प्रथम-क्रम !!{sentence}s~$!E$
यांची संधी $!D_{\Struct{M}}$ असू द्या (ही संधी सांत आहे), म्हणजे
$\Struct{N} \models !D_{\Struct{M}}$ तेव्हा आणि केवळ तेव्हाच, जेव्हा
$\Struct{N} \equiv_n \Struct{M}$. मग $!D = \textstyle\bigvee
\Setabs{!D_{\Struct{M}}}{\Struct{M} \models_L !E}$ असे ठेवा; हा विकल्पही
(तार्किक सममूल्यतेपर्यंत) सांत आहे.

यावरून $\Mod(L){!E} = \Mod(L){!D}$ हा निष्कर्ष मिळतो. वस्तुतः,
$\Struct{N} \models_L !D$ असेल, तर एखाद्या $\Struct{M} \models_L
!E$ साठी $\Struct{N} \models !D_{\Struct{M}}$ असते; त्यामुळे
लेम्माच्या गृहीतकानुसार $\Struct{N} \models_L !E$ हेही खरे असते.
उलट, $\Struct{N} \models_L !E$ असेल, तर $!D_\Struct{N}$ हा $!D$ मधील
एक विकल्पांश आहे; आणि $\Struct{N} \models !D_\Struct{N}$ असल्यामुळे
$\Struct{N} \models_L !D$ हेही खरे असते.
\end{proof}

\begin{thm}[लिंडस्ट्रॉमचे प्रमेय]
  \ollabel{thm:lindstrom} सामान्य अमूर्त तर्कशास्त्र
  $\tuple{L, \models_L}$ ला संहतता आणि लोव्हेनहाइम--स्कोलेम गुणधर्म
  आहेत असे समजा. मग $\tuple{L, \models_L} \le \tuple{F, \models}$
  (म्हणून $\tuple{L, \models_L}$ हे प्रथम-क्रम तर्कशास्त्राशी सममूल्य
  आहे).
  %READERNOTE{OLFOL-049 — गोठवलेल्या इंग्रजीतील प्रमेय-विधानात normal हे आवश्यक गृहीतक राहून गेले आहे; सिद्धता समरूपता, विस्तार आणि सापेक्षीकरण यांसह सामान्यतेचे गुणधर्म स्पष्टपणे वापरते. मराठी विधानात सामान्य हे गृहीतक पुनःस्थापित केले आहे.}
\end{thm}

\begin{proof}
\olref{lem:lindstrom} नुसार, सांत $\Lang{L}$ साठी कोणतेही $!E
\in L(\Lang{L})$ दिले असता असा $n \in \Nat$ असतो की कोणत्याही दोन
!!{structure}s $\Struct{M}$ आणि~$\Struct{N}$ साठी: $\Struct{M}
\equiv_n \Struct{N}$ असेल, तर $\Struct{M}$ आणि $\Struct{N}$ या $!E$
बाबत सहमत असतात, हे दाखवणे पुरेसे आहे. कारण मग $!E$ हे प्रथम-क्रम
!!{sentence} शी सममूल्य असते, आणि त्यावरून $\tuple{L, \models_L} \le
\tuple{F, \models}$ मिळते. आपण सांत आणि निव्वळ संबंधात्मक
!!{language} मध्ये काम करत असल्यामुळे, \olref[bas][pis]{thm:b-n-f}
नुसार $\Struct{M} \equiv_n \Struct{N}$ या विधानाच्या जागी अनुरूप बैजिक
विधान $I_n(\emptyset,\emptyset)$ घेता येते.

$!E$ दिले असता, विरोधासाठी असे समजा की प्रत्येक $n$ साठी अशा
!!{structure}s $\Struct{M}_n$ आणि $\Struct{N}_n$ असतात की
$I_n(\emptyset, \emptyset)$; पण (समजा) $\Struct{M}_n \models_L !E$
आणि $\Struct{N}_n \not\models_L !E$. समरूपता गुणधर्मामुळे सर्व
$\Struct{M}_n$ भाषेतील स्थिरांना त्याच वस्तू नेमतात आणि त्या वस्तूंखेरीज
त्यांचे प्रांत परस्पर वियुक्त आहेत असे आपण मानू शकतो. शिवाय, भाषेत सांत
इतकीच आण्विक !!{sentence}s असल्यामुळे त्या सर्व त्याच आण्विक
!!{sentence}s ची पूर्ति करतात असेही आपण मानू शकतो (नसल्यास
$\Struct{M}_n$ चा उपअनुक्रम घेता येतो). सर्व $\Struct{M}_n$ चा संयोग,
म्हणजे प्रत्येक $\Struct{M}_n$ उपरचना असलेली एकमेव न्यूनतम
!!{structure}, $\Struct{M}$ असू द्या. \olref[lsp]{thm:abstract-p-isom}
च्या सिद्धतेप्रमाणे, $\Struct{M}^*$ ही $\Struct{M}$ ची अशी वर्धित रचना
असू द्या की तिचा !!{domain} $\Domain{M} \cup
\Domain{M}^{<\omega}$ आहे आणि तिच्या विस्तारित !!{language} मध्ये
जोडणी-विधेये $P$ आणि~$Q$ आहेत.
%READERNOTE{OLFOL-050 — सर्व M_n चा रचनात्मक संयोग एकमेव असण्यासाठी त्यांच्या प्रांतांचे छेद स्थिरांच्या समान अर्थनिर्धारणांपुरतेच असणे आवश्यक आहे. गोठवलेल्या इंग्रजीत ही समरूपी प्रती निवडण्याची नेहमीची पायरी न सांगता केवळ स्थिर समान असल्याचे म्हटले आहे; तसेच आवश्यक उपअनुक्रमाला M_n ऐवजी M चा उपअनुक्रम म्हटले आहे. मराठी सिद्धतेत वियुक्तता-अट आणि योग्य n-अधोलिखित दोन्ही पुनःस्थापित केली आहेत.}

$\Struct{N}_n$, $\Struct{N}$ आणि $\Struct{N}^*$ यांची व्याख्या याच
रीतीने करा. आता $\Struct{A}$ ही अशी !!{structure} असू द्या की तिचा
!!{domain} हा $\Struct{M}^*$ आणि $\Struct{N}^*$ यांच्या !!{domain}s
बरोबर नैसर्गिक संख्या~$\Nat$ आणि त्यांचा नैसर्गिक क्रम~$\le$ यांचा
समावेश करतो. तिच्या !!{language} मध्ये !!{domain}s $\Domain{M}$,
$\Domain{N}$, $\Domain{M}^{<\omega}$ आणि $\Domain{N}^{<\omega}$
दर्शवणारी अतिरिक्त विधेये असतात; तसेच $\Struct{M}_n$ आणि
$\Struct{N}_n$ यांचे प्रांत पुढील अर्थाने संकेतबद्ध करणारी विधेये असतात:
\begin{align*}
  \Domain{M_n} & = \Setabs{a \in \Domain{M}}{R(a, n)}; &
  \Domain{N_n} & = \Setabs{a \in \Domain{N}}{S(a,n)}; \\
  \Domain{M}^{<\omega}_n & = \Setabs{a \in \Domain{M}^{<\omega}}{R(a,n)}; &
  \Domain{N}^{<\omega}_n & = \Setabs{a \in \Domain{N}^{<\omega}}{S(a,n)}.
\end{align*}
!!{structure}~$\Struct{A}$ मध्ये असा त्रिस्थानी संबंध $J$ ही असतो की
$J(n, \mathbf{a}, \mathbf{b})$ तेव्हा आणि केवळ तेव्हाच खरे असते,
जेव्हा $I_n(\mathbf{a}, \mathbf{b})$.
%READERNOTE{OLFOL-051 — गोठवलेल्या इंग्रजीत या मोठ्या संकेतबद्ध रचनेसाठी पुन्हा M हेच चिन्ह वापरले आहे, जरी M आधीच सर्व M_n च्या संयोगासाठी वापरलेले आहे. मराठीत मोठ्या रचनेसाठी A वापरून दोन वेगळ्या वस्तू स्पष्ट केल्या आहेत.}

आता $R$, $S$, $J$ इत्यादींनी वाढवलेल्या !!{language}~$\Lang{L}$ मध्ये
असे !!a{sentence}~$!D$ असते की ते $\le$ हा पहिला घटक असलेला पण शेवटचा
घटक नसलेला विविक्त रेषीय क्रम आहे, $\Struct{M}_n \models_L !E$,
$\Struct{N}_n \not\models_L !E$, आणि क्रमातील प्रत्येक $n$ साठी
$J(n, \mathbf{a}, \mathbf{b})$ तेव्हा आणि केवळ तेव्हाच खरे असते,
जेव्हा $I_n(\mathbf{a}, \mathbf{b})$, असे सांगते.
%READERNOTE{OLFOL-052 — E हे अमूर्त तर्कशास्त्रातील वाक्य आहे; म्हणून M_n आणि N_n वरील दोन पूर्ति-चिन्हांना L अधोलिखित आवश्यक आहे. गोठवलेल्या इंग्रजीतील साधे models चिन्ह मराठीत models_L असे दुरुस्त केले आहे.}

संहतता गुणधर्म वापरून, $!D$ बरोबर एका नवीन स्थिराला प्रत्येक मानक
घटकाच्या वर ठेवणारी वाक्ये घ्या. या संचाला असे प्रतिमान $\Struct{A}^*$
मिळते की त्यातील क्रमात अमानक घटक~$n^*$ असतो. विशेषतः,
$\Struct{A}^*$ मध्ये अशा उप!!{structure}s $\Struct{M_{n^*}}$ आणि
$\Struct{N_{n^*}}$ असतात की $\Struct{M_{n^*}} \models_L !E$ आणि
$\Struct{N_{n^*}} \not\models_L !E$. आता $\Domain{M_{n^*}}$ मधील
आणि $\Domain{N_{n^*}}$ मधील $k$-सदस्यी क्रमितकांच्या जोड्यांचा संच
$\mathcal{I}$ पुढीलप्रमाणे परिभाषित करता येतो: $\tuple{\mathbf{a},
\mathbf{b}} \in \mathcal{I}$ तेव्हा आणि केवळ तेव्हाच, जेव्हा
$J(n^*-k, \mathbf{a}, \mathbf{b})$; येथे $k$ ही $\mathbf{a}$ आणि
$\mathbf{b}$ यांची लांबी आहे. $n^*$ हा अमानक असल्यामुळे प्रत्येक मानक
$k$ साठी $n^* - k >0$ असते, आणि $\mathcal{I}$ हा संच
$\Struct{M_{n^*}} \simeq_p \Struct{N_{n^*}}$ या वस्तुस्थितीचा साक्षी
आहे. पण \olref[lsp]{thm:abstract-p-isom} नुसार $\Struct{M_{n^*}}$ ही
$\Struct{N_{n^*}}$ शी $L$-सममूल्य आहे; हा विरोधाभास आहे.
%READERNOTE{OLFOL-053 — एकट्या D ला प्रतिमान असणे क्रमात अमानक घटक सक्तीने देत नाही. नेहमीच्या संहतता-युक्तिवादात नवीन स्थिर प्रत्येक मानक घटकाच्या वर आहे अशी अनंत वाक्येही जोडावी लागतात. गोठवलेल्या इंग्रजीत राहून गेलेली ही वाक्ये मराठी सिद्धतेत स्पष्ट केली आहेत.}
%READERNOTE{OLFOL-054 — गोठवलेल्या इंग्रजीत अंतिम मोठ्या प्रतिमानासाठी M^* या दोन विसंगत TeX रूपांचा वापर आहे आणि ते आधीच्या M^* वर्धित रचनेशीही धडकते. मराठीत संकेतबद्ध मोठ्या प्रतिमानाचे A^* हेच चिन्ह सातत्याने वापरले आहे.}
\end{proof}

\end{document}
